\section{Análisis vertical}

El análisis vertical expresa cada partida como porcentaje de una base común dentro del mismo estado financiero. Permite estudiar la estructura económica o financiera de la empresa.

La formula es:

\[
\text{Participación porcentual} =
\frac{\text{Partida}}{\text{Base}} \times 100
\]

En el Estado de Situación Financiera, la base usual es el total activo. En el Estado de Resultados, la base usual son las ventas netas.

Este método responde a la pregunta: \textbf{¿qué peso tiene una partida dentro del total?}

\subsection{Ejemplo básico de análisis vertical}

Enunciado: una empresa desea conocer qué proporción de sus ventas se destina al costo de ventas.

Datos:

\begin{itemize}
  \item Ventas netas: S/ 100,000.00.
  \item Costo de ventas: S/ 60,000.00.
\end{itemize}

Cálculo:

\[
\frac{\text{S/ 60,000.00}}{\text{S/ 100,000.00}} \times 100 = 60.00\%
\]

Interpretación: de cada S/ 100.00 vendidos, S/ 60.00 se destinan al costo de ventas y S/ 40.00 quedan como margen bruto antes de gastos.

\subsection{Ejemplo aplicado de análisis vertical}

Enunciado: una empresa tecnológica desea evaluar qué peso tienen sus costos de servidores y licencias dentro de sus ventas.

Datos:

\begin{itemize}
  \item Ventas: S/ 200,000.00.
  \item Servidores y licencias: S/ 70,000.00.
\end{itemize}

Cálculo:

\[
\frac{\text{S/ 70,000.00}}{\text{S/ 200,000.00}} \times 100 = 35.00\%
\]

Interpretación y decisión: el 35.00\% de las ventas se destina a infraestructura tecnológica. Si este porcentaje aumenta, la empresa debe revisar arquitectura, proveedores cloud, licencias y escalabilidad.
